(1) 理想四能级激光的有效上能级反转数为 \(N\ge0\)，腔光子数 \(S\ge0\)。忽略自发辐射进入该特定腔模的比例，给定速率方程
\(\dot N=R-\gamma N-gNS\)，\(\dot S=gNS-\kappa S\)，其中 \(R,\gamma,g,\kappa>0\)。找出无激光及激光稳态并求激光起振阈值，解释为何无噪声的精确 \(S(0)=0\) 特殊初值不能自行起振。
(2) 对阈值下和阈值上稳态分别线性化，求线性扰动增长率。判断非激光平衡何时失稳，求激光平衡的稳定性以及在何种参数下存在阻尼松弛振荡。
(3) 对阈值略上的泵浦 \(R=R_{\rm th}(1+\epsilon)\)，\(\epsilon>0\) 较小，求稳态反转钳制、输出光子流率与两种小扰动模的主阶行为；区分直接光子损耗时间 \(\kappa^{-1}\) 与复合系统的临界慢弛豫。
(4) 假设每次有效泵浦增加一个上能级激发，其激光跃迁能量为 \(\hbar\omega_L\)。证明给定方程满足严格的激发能量账本，求稳态腔输出功率、上能级非腔损失和抽运供能；解释实际泵浦光子能量 \(\hbar\omega_p\) 大于 \(\hbar\omega_L\) 时的量子缺陷热。
